% Compile the document with pdflatex, bibliography sequence:
% 1. pdflatex main-bevington-vvuq.tex      % First pass
% 2. biber main-bevington-vvuq             % Compile the bibliography
% 3. pdflatex main-bevington-vvuq.tex      % Second pass
% 4. pdflatex main-bevington-vvuq.tex      % Final pass to resolve references
\documentclass{beamer}
\newcommand{\yshift}{0.0} % Vertical offset for decorative logos

% Fetch home directory to make this file independent of the file system
\usepackage{catchfile}
\CatchFileDef{\HomePath}{|kpsewhich -var-value=HOME}{}
\makeatletter
\edef\HomePath{\expandafter\zap@space\HomePath \@empty}
\makeatother

% Define GitHub and repository paths
\newcommand{\pGithub}{\HomePath/GitHub/}
\newcommand{\pRepo}{\pGithub/latex-2025/}

% Navigation variables: FLAVOR, PROJECT TITLE
\newcommand{\flavor}{pres}        % Type: pres (presentation) or doc (document)
\newcommand{\projTitle}{bevington-vvuq}   % Project Title

\newcommand{\pBibs}{\pRepo/bibs/}
\newcommand{\pGlobal}{\pRepo/global/}
\newcommand{\pConfigRoot}{\pRepo/global/config/}
\newcommand{\pConfigGlobal}{\pConfigRoot/common/}
\newcommand{\pConfigFlavor}{\pConfigRoot/\flavor/}

\newcommand{\enumMain}  { ** MAIN CONFIGURATION ** }
\newcommand{\enumBase}  { == BASE CONFIGURATION == }
\newcommand{\enumFlavor}  { ** FLAVOR CONFIGURATION ** }
\newcommand{\enumLocal} { -- LOCAL CONFIGURATION -- }
\newcommand{\enumAux} { >> AUXILIARY LIBRARY << }
\newcommand{\enumHL} { >> HYPERKLINKS IBRARY << }
\newcommand{\enumGlossary} { >> GLOSSARY LIBRARY << }
\newcommand{\enumElements} { >> ELEMENTS LIBRARY << }
\newcommand{\enumPres}  { ++ PRESENTATION SETUP ++ }
\newcommand{\enumDocs}  { @@ DOCUMENT SETUP @@ }
\newcommand{\enumContent} { !! CONTENT SETUP !! }

% Load configuration files hierarchically
\input{\pConfigGlobal "config-base.tex"}
\input{\pConfigFlavor "config-flavor-dsg.tex"}
\input{\pConfig "config-local.tex"}

\typeout{\enumDocs}
% Choose hyperlink configuration:
\input{\pConfigGlobal/"href-hidden.tex"}   % For hidden links (clean, professional)
% \input{\pConfigGlobal/"href-visible.tex"} % For visible links (debugging, drafts)

%   --   --   --   --   --   --   --   --   --   -- Bibliography format, content
\input{\pConfigGlobal/"bib-config-a.tex"}
\addbibresource{\pBibs/textbooks.bib}

\renewcommand{\listingFontSize}{\scriptsize}

% Document metadata
%   --   --   --   --   --   --   --   --   --   -- Title, Author
\title[VVUQ and the Bevington Example]{A VVUQ Analysis of a \\ Classic Least Squares Example}
\author[\TopaHII]{\TopaDSG \\ \TopaDSGemail}
\institute{\dsgit}
\date{\today}

\begin{document}

% Title slide
\begin{frame}
    \titlepage
\end{frame}

% Welcome audience, frame the problem
\begin{frame}\frametitle{Heat Transfer Via Least Squares}
\begin{center}
	%\href{}{
	\begin{overpic}[ width = \textwidth ]
		{\pGraphics/bevington-temperature-bar-color.pdf}
		%\put(50, 50)	{\colorbox{white}{$a+b$}}
	\end{overpic}
\end{center}
\end{frame}
%
			\input{\pGlobalSlides/"slide-bev-before-after.tex"}
%
\begin{frame}\frametitle{Goal}
\center
	Using this physics-centric problem we will show how \\
	VVUQ can turn a collection of codes into a 
	\bl{robust, reliable scientific tool}.
\end{frame}
%
\begin{frame}\frametitle{Goal}
\center
	A discourse on \href{https://en.wikipedia.org/wiki/Software_verification_and_validation}{\bl{Verification}, \bl{Validation}} and \bl{\uq} using a common example.
\twodots
\end{frame}
%

\begin{frame}[ allowframebreaks ]\frametitle{Presentation Overview}
    \tableofcontents
\end{frame}

%   --   --   --   --   --   --   --   --   --   -- Sections
% Primary content
	\input{\pSections/"sec-introduction.tex"}
	\input{\pSections/"sec-sample-analysis.tex"}
	%\input{\pSections/"sec-code.tex"}
	\input{\pSections/"sec-snapshot.tex"}
	\input{\pSections/"sec-data-structure.tex"}
	%\input{\pSections "sec-backup.tex"}

\appendix
	\input{\pSections/"app-theory.tex"}
	%\input{\pSections "app-normal-eqns.tex"}
	
	%   --   --   --   --   --   --   --   --   --   -- References
\section{} % Do not add to navigation panel
\begin{frame}[ allowframebreaks ]
    \frametitle{References}
    \printbibliography
\end{frame}

% Questions?
\begin{frame}
    \titlepage
\end{frame}

\end{document}

\tiny
\scriptsize
\footnotesize
\small
\normalsize
\large
\Large
\LARGE
\huge
\Huge

% \, thin space (normally 1/6 of a quad);
% \> medium space (normally 2/9 of a quad);
% \; thick space (normally 5/18 of a quad);

\begin{frame}\frametitle{Frame Title}
\begin{enumerate}
	\item First item
	\item Second item
	\item Third item
\end{enumerate}
\end{frame}

\begin{frame}\frametitle{Frame Title}
\begin{equation}
	\begin{array}{ccc} 
			% Example matrix layout
			a & b & c \
			d & e & f \
			g & h & i
	\end{array}
% \label{eq:example}
\end{equation}
\end{frame}

\begin{frame}\frametitle{Graphic Example}
\center
	\href{https://example.com}{
	\begin{overpic}[scale=1.0]{\pLocalGraphics/graphic-file.pdf}
		%\put(-7,-10){Auxiliary text.}
	\end{overpic}}
\end{frame}

\begin{frame}\frametitle{Frame Title}
\begin{table}[htp]
%\caption{Default caption}
\begin{center}
	\begin{tabular}{cc}
		\hline
		Column 1 & Column 2 \
		\hline
		Data 1 & Data 2 \
		Data 3 & Data 4 \
		\hline
	\end{tabular}
\end{center}
% \label{tab:example}
\end{table}
\end{frame}
